% Generated draft: requires mathematical and source review.
\documentclass[11pt]{article}
\usepackage[margin=25mm]{geometry}
\usepackage{fontspec}
\setmainfont{Latin Modern Roman}
\usepackage{amsmath,amssymb,mathtools,unicode-math}
\setmathfont{Latin Modern Math}
\usepackage{enumitem,needspace,xurl,hyperref}
\setlength{\parindent}{0pt}
\setlength{\parskip}{5pt}
\setlength{\emergencystretch}{3em}
\newcommand{\N}{\mathbb{N}}
\newcommand{\Z}{\mathbb{Z}}
\newcommand{\Q}{\mathbb{Q}}
\newcommand{\R}{\mathbb{R}}
\newcommand{\C}{\mathbb{C}}
\newcommand{\F}{\mathbb{F}}
\newcommand{\A}{\mathbb{A}}
\newcommand{\PP}{\mathbb{P}}
\newcommand{\E}{\mathbb{E}}
\newcommand{\eps}{\varepsilon}
\newcommand{\dd}{\,\mathrm{d}}
\begin{document}
\section*{UnsolvedMath ID: 30005608}
\textbf{Source problem code:} catalogue entry
\subsection{Definitions and mathematical statement}
For notation, let \(\Omega\) be a bounded domain in a Riemannian manifold \((M,g)\), with boundary regularity sufficient to define a discrete Robin spectrum. Under the convention
\[
-\Delta_g u=\lambda u\quad\text{in }\Omega,
\qquad
\partial_\nu u+\alpha u=0\quad\text{on }\partial\Omega,
\]
where \(\nu\) is the outward unit normal and \(\alpha\in\R\), write
\[
\lambda_1(\Omega;\alpha)\leq\lambda_2(\Omega;\alpha)\leq\cdots
\]
for the eigenvalues counted with multiplicity. These conventions provide notation; the supplied record does not specify the source's Robin sign convention or admissible domain class. An eigenfunction on a rotationally symmetric comparison cap is called radial if it depends only on distance from the cap's center.

\textbf{Requested maximality problem, with unresolved qualifications.} The supplied reconstructed statement asks for spherical-cap maximality of the second Robin eigenvalue for hyperbolic domains with positive Robin parameter and in the regime where the cap's second eigenfunction is radial. If the intended comparison holds the Robin parameter fixed, its inequality would have the form
\[
\lambda_2(\Omega;\alpha)\leq\lambda_2(C;\alpha)
\qquad\text{for every admissible comparison domain }\Omega,
\]
where \(C\) denotes the intended comparison cap. This is only a conditional template: the record does not identify the ambient dimension, curvature normalization, meaning of the cap in the hyperbolic setting, constraint defining the comparison class, boundary regularity, connectedness assumptions, or whether the Robin parameter itself is fixed or normalized with domain size. Consequently, the quantifier over admissible domains cannot yet be made precise.

There is also an unresolved discrepancy in the target parameter set. The retained original passage says that the method is not applicable for hyperbolic domains with positive Robin parameter, described as the ``2nd quadrant,'' \emph{or} where the cap's second eigenfunction is radial, described as the ``red region.'' The reconstructed statement instead joins the two descriptions with ``and.'' The record does not establish whether these are two separate regimes, their intersection, or regimes in different ambient geometries. It supplies neither the referenced diagram nor its parameter definitions. Moreover, if the second eigenvalue is multiple, the phrase ``the second eigenfunction is radial'' does not specify whether a radial eigenfunction merely exists in that eigenspace or whether a particular spectral branch is intended. These qualifications must be resolved before this template constitutes a fully specified conjecture. Failure of the cited method alone does not establish the maximality inequality.
\subsection{Short English statement}
The supplied problem asks whether a comparison cap maximizes the second Robin eigenvalue in parameter regimes beyond the cited method. A precise conjecture remains unresolved because the comparison constraint and geometric setting are unspecified, and the original passage uses ``or'' where the reconstructed statement uses ``and.'' The record labels the problem open; this draft does not independently verify that status.
\subsection{Sources}
\begin{itemize}
\item Local UnsolvedMath registry: \texttt{data/unsolvedmath/problems.json}, with \texttt{catalog/problems.json} as the additional-record registry; supplied permanent ID \(30005608\), title \emph{Spherical Robin-Eigenvalue Maximization Beyond Known Regimes}, exact source problem code \texttt{catalogue entry}. This draft uses the supplied complete record; registry files were not independently inspected.
\item Supplied report citation: \emph{Geometric Spectral Theory (2024)}, Oberwolfach Reports, pp. 2017--2106, DOI \texttt{10.4171/owr/2023/36}. Supplied source URL and literature reference labeled ``Oberwolfach Report 36'': \url{https://doi.org/10.4171/owr/2023/36}. The record attributes the problem to an Oberwolfach workshop report and gives proposed year 2023, with no individual proposer identified.
\item Supplied literature assessment, checked \textbf{2026-08-22}: status ``open.'' The assessment says that the source treats the maximality question in regions outside the existing method as open and that no general replacement argument was found. This is supplied dated curation, not an independently verified literature assessment; the record expressly requests independent verification of time-sensitive status and cited claims.
\item Supplied statement metadata: status \texttt{corrected\_verified}, with the assertion that the reconstruction was checked against the official Oberwolfach report. This assertion is preserved as source metadata and does not constitute verification by this draft.
\end{itemize}
\end{document}
